\chapter{主成分分析}
本书四章按研究目的可分为两类：主成分分析与因子分析以\textbf{变量}为对象，目的是将多个相关指标概括为少数几个综合指标；聚类分析与判别分析以\textbf{观察对象}为对象，目的是将个体（或样品）分类，或判定新个体所属的类别。

主成分分析是后续各章的基础。以6项智力测验得分为例，各项得分彼此相关，所含信息有重叠，主成分分析的任务是用一两个综合得分代替原有指标，并尽可能少地损失信息。这里的“信息”有明确的定义，即综合得分在个体间的方差，方差越大，区分个体的能力越强。本章的数学推导即求方差最大的线性组合，其解为协方差阵（或相关阵）的特征向量。

主成分分析属于数据变换与描述方法，不涉及统计模型，也不要求分布假设。它不假定数据背后存在潜在因子，这是它与第2章因子分析的根本区别。
\section{主成分分析}
主成分分析（principal component analysis，PCA）是利用多个数值变量的相关性，把它们变换成少数几个互不相关的综合指标（主成分），以反映多个变量所提供的信息。
\Note{PCA靠变量间的相关性来降维。若标准化后的变量近似互不相关，即相关阵接近单位阵，各特征值都接近1，没有少数几个主成分能概括大部分方差，降维效果就差。用协方差阵分析时，即使变量互不相关，方差特别大的变量也会单独占据前几个主成分，但这只反映量纲或尺度的差异，并没有实现降维。}
\Example{例1\quad 六项智力测量得分}
寻找一个或几个综合指标，概括六项得分的信息。
\par\addvspace{5pt}\begin{center}\begin{minipage}{\linewidth}\centering
\small
\captionof{table}{某小学10名男学生六个项目的智力测量计分}\label{tab:iq}
\begin{tabular}{lrrrrrr}\toprule
{\bfseries 编号 }&{\bfseries  常识 $X_1$ }&{\bfseries  算术 $X_2$ }&{\bfseries  理解 $X_3$ }&{\bfseries  填图 $X_4$ }&{\bfseries  积木 $X_5$ }&{\bfseries  译码 $X_6$}\\\midrule
1 & 14 & 13 & 28 & 14 & 22 & 39\\
2 & 10 & 14 & 15 & 14 & 34 & 35\\
3 & 11 & 12 & 19 & 13 & 24 & 39\\
4 & 7 & 7 & 7 & 9 & 20 & 23\\
5 & 13 & 12 & 24 & 12 & 26 & 38\\
6 & 19 & 14 & 22 & 16 & 23 & 37\\
7 & 20 & 16 & 26 & 21 & 38 & 69\\
8 & 9 & 10 & 14 & 9 & 31 & 46\\
9 & 9 & 8 & 15 & 13 & 14 & 46\\
10 & 9 & 9 & 12 & 10 & 23 & 46\\
\bottomrule\end{tabular}
\end{minipage}\end{center}

\Source{资料来源：余松林，《医学现场研究中的统计分析方法》，同济医科大学（内部教材），143页。}
\subsection{六项得分的协方差阵和相关矩阵}
\begin{code}
\RComment{\# 按表\ref{tab:iq}录入数据}
d \ROperator{<-} \RFunction{data.frame}(
  x1 \ROperator{=} \RFunction{c}(\RNumber{14, 10, 11, 7, 13, 19, 20, 9, 9, 9}),
  x2 \ROperator{=} \RFunction{c}(\RNumber{13, 14, 12, 7, 12, 14, 16, 10, 8, 9}),
  x3 \ROperator{=} \RFunction{c}(\RNumber{28, 15, 19, 7, 24, 22, 26, 14, 15, 12}),
  x4 \ROperator{=} \RFunction{c}(\RNumber{14, 14, 13, 9, 12, 16, 21, 9, 13, 10}),
  x5 \ROperator{=} \RFunction{c}(\RNumber{22, 34, 24, 20, 26, 23, 38, 31, 14, 23}),
  x6 \ROperator{=} \RFunction{c}(\RNumber{39, 35, 39, 23, 38, 37, 69, 46, 46, 46}))
S \ROperator{<-} \RFunction{cov}(d)
S
\RFunction{cor}(d)
\RFunction{diag}(S)
\RFunction{sum}(\RFunction{diag}(S))
\RFunction{diag}(\RFunction{cor}(d))
\end{code}
\par\addvspace{5pt}\begin{center}\begin{minipage}{\linewidth}\centering
\small
\captionof{table}{协方差阵 $S$}
\begin{tabular}{lrrrrrr}\toprule
 &{\bfseries  $X_1$ }&{\bfseries  $X_2$ }&{\bfseries  $X_3$ }&{\bfseries  $X_4$ }&{\bfseries  $X_5$ }&{\bfseries  $X_6$}\\\midrule
$X_1$ & 19.433333 & 10.722222 & 24.088889 & 13.877778 & 12.611111 & 27.466667\\
$X_2$ & 10.722222 & 8.500000 & 15.333333 & 8.722222 & 14.277778 & 15.444444\\
$X_3$ & 24.088889 & 15.333333 & 45.288889 & 17.200000 & 13.222222 & 35.266667\\
$X_4$ & 13.877778 & 8.722222 & 17.200000 & 12.988889 & 11.611111 & 27.022222\\
$X_5$ & 12.611111 & 14.277778 & 13.222222 & 11.611111 & 49.833333 & 41.555556\\
$X_6$ & 27.466667 & 15.444444 & 35.266667 & 27.022222 & 41.555556 & 138.400000\\
\bottomrule\end{tabular}
\end{minipage}\end{center}

\par\addvspace{5pt}\begin{center}\begin{minipage}{\linewidth}\centering
\small
\captionof{table}{相关矩阵 $R$}
\begin{tabular}{lrrrrrr}\toprule
 &{\bfseries  $X_1$ }&{\bfseries  $X_2$ }&{\bfseries  $X_3$ }&{\bfseries  $X_4$ }&{\bfseries  $X_5$ }&{\bfseries  $X_6$}\\\midrule
$X_1$ & 1.0000000 & 0.8342605 & 0.8119837 & 0.8734947 & 0.4052469 & 0.5296199\\
$X_2$ & 0.8342605 & 1.0000000 & 0.7815041 & 0.8301024 & 0.6937311 & 0.4502926\\
$X_3$ & 0.8119837 & 0.7815041 & 1.0000000 & 0.7091639 & 0.2783227 & 0.4454516\\
$X_4$ & 0.8734947 & 0.8301024 & 0.7091639 & 1.0000000 & 0.4563813 & 0.6373339\\
$X_5$ & 0.4052469 & 0.6937311 & 0.2783227 & 0.4563813 & 1.0000000 & 0.5003813\\
$X_6$ & 0.5296199 & 0.4502926 & 0.4454516 & 0.6373339 & 0.5003813 & 1.0000000\\
\bottomrule\end{tabular}
\end{minipage}\end{center}

\[
\diag(S)=(19.43333,\ 8.50000,\ 45.28889,\ 12.98889,\ 49.83333,\ 138.40000),
\]
\[
\tr(S)=274.4444,\qquad \diag(R)=(1,1,1,1,1,1).
\]
实际分析一般调用\rinline{prcomp()}。下面先列出结果，各数值的来源见后续各节。\rinline{prcomp()}默认基于协方差阵，设定\rinline{scale.=TRUE}时基于相关阵：
\par\noindent\begin{rcode}
pc <- prcomp(d)                 # 协方差阵
summary(pc)
round(pc$rotation[,1:2],3)      # 前两个特征向量
\end{rcode}
\begin{routput}
Importance of components:
                           PC1    PC2    PC3     PC4     PC5     PC6
Standard deviation     13.6558 6.6942 5.8888 2.53505 1.31898 0.55498
Proportion of Variance  0.6795 0.1633 0.1264 0.02342 0.00634 0.00112
Cumulative Proportion   0.6795 0.8428 0.9691 0.99254 0.99888 1.00000
      PC1    PC2
x1 -0.232  0.377
x2 -0.149  0.261
x3 -0.317  0.710
x4 -0.207  0.210
x5 -0.334  0.021
x6 -0.818 -0.491
\end{routput}
\rinline{Standard deviation}的平方为特征值，如$13.6558^2=186.48$；\rinline{Proportion of Variance}为贡献率，如$186.48/274.44=0.6795$。特征向量整体变号不影响结果（见第\ref{sec:pc-prop}节后的说明），此处R给出的第一特征向量各分量均为负值，解释时可统一乘以$-1$。

第一主成分中译码$x_6$的系数达0.818。六项测验的满分与分布不同，译码的方差（138.4）为算术（8.5）的16倍，因而在协方差阵分析中占主导地位。改用相关阵的结果如下：
\par\noindent\begin{rcode}
pr <- prcomp(d,scale.=TRUE)     # 相关阵
summary(pr)
round(pr$rotation[,1:2],3)
\end{rcode}
\begin{routput}
Importance of components:
                          PC1    PC2    PC3     PC4     PC5    PC6
Standard deviation     2.0364 0.9285 0.7759 0.50682 0.32673 0.1588
Proportion of Variance 0.6912 0.1437 0.1003 0.04281 0.01779 0.0042
Cumulative Proportion  0.6912 0.8348 0.9352 0.97800 0.99580 1.0000
      PC1    PC2
x1 -0.450  0.289
x2 -0.458 -0.005
x3 -0.408  0.448
x4 -0.453  0.119
x5 -0.315 -0.747
x6 -0.341 -0.379
\end{routput}
\begin{sikao}[主成分的解释]
基于相关阵时，第一主成分的六个系数同号且大小接近（0.32--0.46），可解释为总体智力水平，得分高者各项测验成绩均较好。它与六项标准化得分之和的相关系数为$0.9985$（符号取正后），接近于标准化总分。其原因在于：六项指标两两正相关且程度相近时，方差最大的方向接近等权重方向。

第二主成分的系数有正有负：常识（0.289）、理解（0.448）为正，积木（$-0.747$）、译码（$-0.379$）为负。该主成分为一种对比，得分高表示言语类项目相对较好、操作类项目相对较差，反映能力结构的差异，与总体水平无关。

解释主成分时主要依据系数的符号与相对大小。第二主成分及以后的主成分常为此类对比，其专业意义需结合专业知识判断；难以解释时不必强行命名。
\end{sikao}
\section{主成分的计算方法}
设$p$个随机变量$X=(X_1,\ldots,X_p)\trans$的总体均数为$\mu$、协方差阵为$\Sigma$。总体主成分是中心化变量的线性组合$z_i=a_i\trans(X-\mu)$。实际计算时用样本均数$\bar x$和样本协方差阵$S$代替$\mu$和$\Sigma$，得到样本主成分，其方差为$S$的特征值$\widehat\lambda_i$。基于相关阵时，先把各变量标准化。样本形式为
\begin{align}
z_i&=a_{i1}(x_1-\bar x_1)+a_{i2}(x_2-\bar x_2)+\cdots+a_{ip}(x_p-\bar x_p),\\
c_i&=a_{i1}x_1^*+a_{i2}x_2^*+\cdots+a_{ip}x_p^*,\qquad x_j^*=\frac{x_j-\bar x_j}{s_j}.
\end{align}
前者基于协方差阵提取主成分，后者基于相关矩阵提取主成分。$p$个变量最多有$p$个主成分。$n$例样本的中心化数据阵秩不超过$n-1$，故$\operatorname{rank}(S)\leqslant\min(p,n-1)$，方差非零的样本主成分不超过$\min(p,n-1)$个。
\Note{要让主成分的方差最大，才更有代表性。}
\subsection{变量线性组合的性质}
设$p$个随机变量$X$的均数为$\mu$，协方差阵为$\Sigma$。$p$个变量的线性组合$a\trans X$的均数为$a\trans\mu$，方差为$a\trans\Sigma a$。
\[
X=\begin{bmatrix}X_1\\\vdots\\X_p\end{bmatrix},\qquad
 a=\begin{bmatrix}a_1\\\vdots\\a_p\end{bmatrix},\qquad
 a\trans X=\begin{bmatrix}a_1&\cdots&a_p\end{bmatrix}X
 =a_1X_1+\cdots+a_pX_p.
\]
\subsection{用总分或平均分表示学生智力}
以总分或平均分作为综合指标，是将各变量等同对待，赋予相同权。
\[
a\trans X=\begin{bmatrix}1&\cdots&1\end{bmatrix}
\begin{bmatrix}X_1\\\vdots\\X_p\end{bmatrix}=X_1+\cdots+X_p,
\]
\[
a\trans X=\begin{bmatrix}1/6&\cdots&1/6\end{bmatrix}
\begin{bmatrix}X_1\\\vdots\\X_p\end{bmatrix}=\frac16X_1+\cdots+\frac16X_p.
\]
\Note{这样处理，还不能尽可能多地概括全部变量的信息。}
下面比较等权组合与第一主成分的方差。为使比较对等，将等权重也标准化为单位长度，$a=(1/\sqrt6,\ldots,1/\sqrt6)\trans$：
\par\noindent\begin{rcode}
a <- rep(1/sqrt(6),6)
c(t(a)%*%cov(d)%*%a, pc$sdev[1]^2)   # 协方差阵：等权组合与第一主成分的方差
c(t(a)%*%cor(d)%*%a, pr$sdev[1]^2)   # 相关阵
\end{rcode}
\begin{routput}
[1] 141.8815 186.4801
[1] 4.079090 4.146965
\end{routput}
按原始分数等权相加，方差为141.9，第一主成分为186.5；按标准化分数等权相加，方差为4.079，与第一主成分的4.147接近。第一主成分是所有单位长度权重中方差最大者，等权组合只是其中之一。相关阵下二者接近，原因如前所述；协方差阵下二者差距较大，因为各项方差不同，最优权重偏向方差较大的项目。
\subsection{基于协方差阵提取第一主成分}
\[
z_1=a_{11}(x_1-\bar x_1)+a_{12}(x_2-\bar x_2)+\cdots+a_{1p}(x_p-\bar x_p).
\]
要求$a_1\trans a_1=1$，且第一主成分$z_1$的方差最大。若不加单位长度约束，把$a_1$乘以常数$c$，方差就变为原来的$c^2$倍，可以任意大。加上约束后，问题只比较组合的方向。

根据瑞利商性质，
\begin{equation}
\lambda_{\min}\leqslant\frac{x\trans Ax}{x\trans x}\leqslant\lambda_{\max}.
\end{equation}
当$A=\Sigma$、$x=e_1$时，式子的值最大，为$\lambda_1$。$e_1$是第一特征向量。
\begin{zm}
由谱分解，$\Sigma=E\Lambda E\trans$，其中$E=(e_1,\ldots,e_p)$为正交阵，$\Lambda=\diag(\lambda_1,\ldots,\lambda_p)$，$\lambda_1\geqslant\cdots\geqslant\lambda_p\geqslant0$。令$b=E\trans a$，则$b\trans b=a\trans a=1$，且
\[
a\trans\Sigma a=b\trans\Lambda b=\sum_{j=1}^p\lambda_jb_j^2\leqslant\lambda_1\sum_{j=1}^pb_j^2=\lambda_1,
\]
当$b=(1,0,\ldots,0)\trans$，即$a=e_1$时取等号。
\end{zm}
\subsection{基于协方差阵提取第二主成分}
如果$z_1$不能很好地概括原指标信息，就继续提取第二主成分：
\[
z_2=a_{21}(x_1-\bar x_1)+a_{22}(x_2-\bar x_2)+\cdots+a_{2p}(x_p-\bar x_p).
\]
要求$a_2\trans a_2=1$、$a_1\trans a_2=0$，且第二主成分$z_2$的方差次大。$e_2$是$z_2$表达式的系数，$\var(z_2)=\lambda_2$。
\begin{zm}
沿用上面的记号。约束$a_2\trans e_1=0$即$b_1=0$，于是$a_2\trans\Sigma a_2=\sum_{j\geqslant2}\lambda_jb_j^2\leqslant\lambda_2\sum_{j\geqslant2}b_j^2=\lambda_2$，在$a_2=e_2$时取等号。依此类推，要求与前$i-1$个系数向量正交时，第$i$个主成分的系数是$e_i$，方差为$\lambda_i$。
\end{zm}
\Note{第二主成分与第一主成分两两不相关（见\ref{sec:pc-prop}节性质2），一般并不独立。最多有$p$个主成分，本例即为6个。6个特征值之和等于协方差阵的迹$\tr(S)=274.44$。原例标注的0.679是第一主成分的贡献率$\widehat\lambda_1/\tr(S)$。}
\subsection{基于相关矩阵计算主成分}
对$p\times p$相关矩阵求特征值$\lambda_i$和特征向量$e_i$，$i=1,2,\ldots,p$。
\[
c_1=a_{11}x_1^*+a_{12}x_2^*+\cdots+a_{1p}x_p^*.
\]
$e_1$是第一主成分$c_1$表达式的系数$a_1$，$\var(c_1)=\lambda_1$最大。$e_2$是第二主成分$c_2$表达式的系数，$\var(c_2)=\lambda_2$次大，以此类推。
\Note{标准化变量的方差都是1，总方差为$\tr(R)=p$。这里的“信息量”指方差。}
\section{主成分的性质}\label{sec:pc-prop}
\begin{enumerate}
\item 主成分$c_i$（$i=1,2,\ldots,p$）的均数为0。
\item 主成分$c_i$和$c_j$（$i\ne j$）两两不相关，$\cov(c_i,c_j)=0$。
\item $\var(c_1)=\lambda_1$最大，$\var(c_2)=\lambda_2$次大，以此类推。
\item 总方差不增不减。协方差阵分析保留原变量的总方差$\tr(\Sigma)$，相关阵分析保留标准化变量的总方差$\tr(R)=p$：
\[
\sum_{i=1}^{p}\var(z_i)=\sum_{i=1}^p\lambda_i=\tr(\Sigma)=\sum_{i=1}^{p}\var(x_i),\qquad
\sum_{i=1}^{p}\var(c_i)=\tr(R)=p.
\]
\item 主成分与原始指标间的相关系数（$u_{i1}$为第一特征向量$e_1$的第$i$个分量）：
\[
\rho_{i1}=\frac{u_{i1}\sqrt{\lambda_1}}{\sigma_i}\quad\text{（基于协方差阵）},\qquad
\rho_{i1}=u_{i1}\sqrt{\lambda_1}\quad\text{（基于相关矩阵）}.
\]
这是因为$\cov(x_i,z_1)=(\Sigma e_1)_i=\lambda_1u_{i1}$，再除以$\sigma_i\sqrt{\lambda_1}$。
\end{enumerate}
\begin{zm}[性质2]
$\cov(z_i,z_j)=\cov(e_i\trans X,e_j\trans X)=e_i\trans\Sigma e_j=\lambda_je_i\trans e_j=0$（$i\ne j$）。基于相关阵时把$\Sigma$换成$R$即可。
\end{zm}
不相关不等于独立。若$X$服从多元正态分布，主成分是$X$的线性变换，也服从联合正态分布，这时不相关与独立等价，总体主成分相互独立。样本主成分得分直接满足样本协方差为零：得分阵$Z=X_cE$的样本协方差阵为$E\trans SE=\widehat\Lambda$，为对角阵。

\paragraph{方差贡献率与累计贡献率}
第$j$个主成分的方差贡献率为$\lambda_j/\sum_{i=1}^p\lambda_i$，前$k$个主成分的累计贡献率为$\sum_{j=1}^k\lambda_j/\sum_{i=1}^p\lambda_i$。协方差阵分析的分母为$\tr(\Sigma)$，相关阵分析的分母为$p$。例1基于协方差阵：前两个主成分的贡献率分别为67.9\%和16.3\%，累计84.3\%。
\Example{例2\quad 29名13岁男童的身高、体重和肺活量}
\par\addvspace{5pt}\begin{center}\begin{minipage}{\linewidth}\centering
\small
\captionof{table}{某地29名13岁男童的身高、体重和肺活量}\label{tab:hwv}
\begin{tabular}{lrrrrrrr}\toprule
{\bfseries 编号 }&{\bfseries  身高/cm }&{\bfseries  体重/kg }&{\bfseries  肺活量/L }&{\bfseries  编号 }&{\bfseries  身高/cm }&{\bfseries  体重/kg }&{\bfseries  肺活量/L}\\\midrule
1 & 135.1 & 32.0 & 1.75 & 2 & 139.9 & 30.4 & 2.00\\
3 & 163.6 & 46.2 & 2.75 & 4 & 146.5 & 33.5 & 2.50\\
5 & 156.2 & 37.1 & 2.75 & 6 & 156.4 & 35.5 & 2.00\\
7 & 167.8 & 41.5 & 2.75 & 8 & 149.7 & 31.0 & 1.50\\
9 & 145.0 & 33.0 & 2.50 & 10 & 148.5 & 37.2 & 2.25\\
11 & 165.5 & 49.5 & 3.00 & 12 & 135.0 & 27.6 & 1.25\\
13 & 153.3 & 41.0 & 2.75 & 14 & 152.0 & 32.0 & 1.75\\
15 & 160.5 & 47.2 & 2.25 & 16 & 153.0 & 32.0 & 1.75\\
17 & 147.6 & 40.5 & 2.00 & 18 & 157.5 & 43.3 & 2.25\\
19 & 155.1 & 44.7 & 2.75 & 20 & 160.5 & 37.5 & 2.00\\
21 & 143.0 & 31.5 & 1.75 & 22 & 149.4 & 33.9 & 2.25\\
23 & 160.8 & 40.4 & 2.75 & 24 & 159.0 & 38.5 & 2.50\\
25 & 158.2 & 37.5 & 2.00 & 26 & 150.0 & 36.0 & 1.75\\
27 & 144.5 & 34.7 & 2.25 & 28 & 154.6 & 39.5 & 2.50\\
29 & 156.5 & 32.0 & 1.75 &  &  &  & \\
\bottomrule\end{tabular}
\end{minipage}\end{center}

\Source{资料来源：陈峰主编，《现代医学统计方法与Stata应用》，中国统计出版社，2000，109页。}
\subsection{协方差阵与相关阵的选择}
协方差阵分析按各变量的方差赋权，方差大的变量在第一主成分中占主导，结果随测量单位改变。相关阵分析先把每个变量标准化，使各变量的初始方差都为1，即在进入分析前赋予相同的方差权重，结果不随单位的线性变换而改变。用表\ref{tab:hwv}的数据复算：
\begin{center}\small\begin{tabular}{lcc}\toprule
{\bfseries 分析方式}&{\bfseries 第一主成分贡献率}&{\bfseries $e_1$（身高，体重，肺活量）}\\\midrule
协方差阵（cm，kg，L）&89.3\%&$(0.865,\ 0.501,\ 0.031)$\\
协方差阵（m，kg，mL）&99.99\%&$(0.000,\ 0.009,\ 1.000)$\\
相关阵（任意单位）&79.3\%&$(0.565,\ 0.604,\ 0.563)$\\\bottomrule
\end{tabular}\end{center}
以cm、kg、L为单位时，身高方差最大，第一主成分主要反映身高；把身高改用m、肺活量改用mL后，第一主成分几乎只剩肺活量。贡献率看似接近100\%，只是因为肺活量的数值方差被放大。相关阵分析在两种单位下完全相同。因此，变量单位不同或方差相差悬殊时宜用相关阵；只有各变量单位相同、方差大小本身有专业意义时（如例4的空气污染分指数），才用协方差阵。
\section{主成分的几何意义}
\Example{例3\quad 18名学生的语文、数学成绩}
从18名学生的语文、数学成绩中，基于协方差阵或相关阵提取主成分，进行综合评价。原始成绩即图\ref{fig:scores}左图中各点的坐标。
\begin{code}
d \ROperator{<-} \RFunction{data.frame}(
  x1 \ROperator{=} \RFunction{c}(\RNumber{45, 48, 52, 61, 65, 67, 73, 75, 78, 81, 84, 88, 89, 93, 95, 98, 99, 98}),
  x2 \ROperator{=} \RFunction{c}(\RNumber{43, 53, 58, 70, 61, 70, 75, 78, 83, 88, 76, 92, 78, 82, 97, 95, 92, 89}))
\RFunction{eigen}(\RFunction{cor}(d))
\RFunction{acos}(\RNumber{0.7071068}) \ROperator{*} \RNumber{180} \ROperator{/} pi
\end{code}
\[
\lambda_1=1.93200236,\quad \lambda_2=0.06799764,\qquad
\begin{bmatrix}e_1&e_2\end{bmatrix}=
\begin{bmatrix}0.7071068&-0.7071068\\0.7071068&0.7071068\end{bmatrix}.
\]
\[
c_1=0.707x_1^*+0.707x_2^*,\qquad
\var(c_1)=\lambda_1=1.932,\qquad \frac{\lambda_1}{p}=\frac{1.932}{2}=96.6\%.
\]
\begin{figure}[!htbp]\centering
\begin{tikzpicture}\begin{axis}[width=.46\textwidth,height=6.8cm,xlabel={$x_1$},ylabel={$x_2$},xmin=42,xmax=101,ymin=40,ymax=100]\addplot[only marks,mark=*,mark size=2.1pt,Brick] coordinates {(45,43) (48,53) (52,58) (61,70) (65,61) (67,70) (73,75) (75,78) (78,83) (81,88) (84,76) (88,92) (89,78) (93,82) (95,97) (98,95) (99,92) (98,89)};\addplot[densely dashed,line width=.45pt,Muted] coordinates {(77.16666666666667,40) (77.16666666666667,100)};\addplot[densely dashed,line width=.45pt,Muted] coordinates {(42,76.66666666666667) (101,76.66666666666667)};\end{axis}\end{tikzpicture}
\hfill
\begin{tikzpicture}\begin{axis}[width=.46\textwidth,height=6.8cm,xlabel={$x_1^*$},ylabel={$x_2^*$},xmin=-2,xmax=1.5,ymin=-2.4,ymax=1.5]\addplot[only marks,mark=*,mark size=2.1pt,Forest] coordinates {(-1.8240906,-2.2080866) (-1.6539681,-1.5522193) (-1.4271382,-1.2242857) (-0.91677091,-0.43724488) (-0.689941,-1.0275255) (-0.57652604,-0.43724488) (-0.23628116,-0.10931122) (-0.12286621,0.087448976) (0.047256233,0.41538264) (0.21737867,0.7433163) (0.38750111,-0.043724488) (0.61433103,1.0056632) (0.6710385,0.087448976) (0.89786842,0.3497959) (1.0112834,1.3335969) (1.1814058,1.2024234) (1.2381133,1.0056632) (1.1814058,0.80890303)};\addplot[densely dashed,line width=.45pt,Muted] coordinates {(0,-2.4) (0,1.5)};\addplot[densely dashed,line width=.45pt,Muted] coordinates {(-2,0) (1.5,0)};\end{axis}\end{tikzpicture}
\caption{语文、数学成绩及其标化成绩}\label{fig:scores}
\end{figure}\FloatBarrier
\begin{figure}[!htbp]\centering
\begin{tikzpicture}\begin{axis}[width=9.5cm,height=8.7cm,xlabel={$x_1^*$},ylabel={$x_2^*$},xmin=-2,xmax=2,ymin=-2.3,ymax=2,axis equal image]\addplot[only marks,mark=*,mark size=2pt,Ink!75] coordinates {(-1.8240906,-2.2080866) (-1.6539681,-1.5522193) (-1.4271382,-1.2242857) (-0.91677091,-0.43724488) (-0.689941,-1.0275255) (-0.57652604,-0.43724488) (-0.23628116,-0.10931122) (-0.12286621,0.087448976) (0.047256233,0.41538264) (0.21737867,0.7433163) (0.38750111,-0.043724488) (0.61433103,1.0056632) (0.6710385,0.087448976) (0.89786842,0.3497959) (1.0112834,1.3335969) (1.1814058,1.2024234) (1.2381133,1.0056632) (1.1814058,0.80890303)};\addplot[densely dotted,line width=.45pt,Muted] coordinates {(-2,0) (2,0)};\addplot[densely dotted,line width=.45pt,Muted] coordinates {(0,-2.3) (0,2)};\addplot[line width=1.1pt,Brick,<->] coordinates {(-2,-2) (2,2)} node[pos=.94,above] {$c_1$};\addplot[line width=1.1pt,Forest,dashed,<->] coordinates {(-2,2) (2,-2)} node[pos=.92,above] {$c_2$};\end{axis}\end{tikzpicture}
\caption{主成分坐标轴旋转}
\end{figure}\FloatBarrier

原坐标系中，语文与数学成绩有相关。基于二维相关阵作主成分分析，将坐标轴逆时针旋转$45^\circ$，形成新坐标系。数据点在$c_1$轴上的投影最分散，即方差最大，在$c_2$轴上的投影相对紧密。$c_1$与$c_2$不相关，$c_1$包含大量原始信息，可作进一步分析。
\subsection{旋转矩阵}
\[
T_{\text{顺时针}}=\begin{bmatrix}\cos\theta&\sin\theta\\-\sin\theta&\cos\theta\end{bmatrix},\qquad
T_{\text{逆时针}}=\begin{bmatrix}\cos\theta&-\sin\theta\\\sin\theta&\cos\theta\end{bmatrix}.
\]
\subsection{保留主成分数与重构误差}
设$X_c$为$n\times p$的中心化（或标准化）数据阵，$S=X_c\trans X_c/(n-1)=E\widehat\Lambda E\trans$，得分阵$Z=X_cE$。只保留前$k$个主成分，记$E_k=(e_1,\ldots,e_k)$，用这$k$个得分重构原数据：
\[
\widehat X_c=Z_kE_k\trans=X_cE_kE_k\trans.
\]
\begin{zm}
记$\widetilde E=(e_{k+1},\ldots,e_p)$。由$EE\trans=I$，$X_c-\widehat X_c=X_c(I-E_kE_k\trans)=X_c\widetilde E\widetilde E\trans$，于是
\[
\|X_c-\widehat X_c\|_F^2=\tr\bigl(\widetilde E\trans X_c\trans X_c\widetilde E\bigr)=(n-1)\tr\bigl(\widetilde E\trans S\widetilde E\bigr)=(n-1)\sum_{j>k}\widehat\lambda_j.
\]
再证这是最小值。任取$p\times k$列正交阵$V$（$V\trans V=I_k$），把数据投影到$V$张成的$k$维子空间，重构误差为$(n-1)\bigl[\tr(S)-\tr(V\trans SV)\bigr]$。令$w_j=\|V\trans e_j\|^2$，则$0\leqslant w_j\leqslant1$，$\sum_jw_j=\tr(V\trans V)=k$，且$\tr(V\trans SV)=\sum_j\widehat\lambda_jw_j\leqslant\widehat\lambda_1+\cdots+\widehat\lambda_k$。$V=E_k$时取等号，故前$k$个主成分给出最小平方重构误差。对一般秩$k$近似阵，同样结论即Eckart--Young定理。
\end{zm}
因此
\[
\frac{\|X_c-\widehat X_c\|_F^2}{\|X_c\|_F^2}=\frac{\sum_{j>k}\widehat\lambda_j}{\sum_{j=1}^p\widehat\lambda_j}=1-\text{前$k$个主成分的累计贡献率}.
\]
累计贡献率就是前$k$个主成分重构出的那部分总方差（平方和）所占比例。碎石图把$\widehat\lambda_j$按$j$作图，曲线由陡变平的“拐点”之后，每多保留一个主成分，重构误差的减少都很有限。常用的经验做法是取累计贡献率达到70\%～85\%、碎石图拐点前的主成分，或在相关阵分析中保留特征值大于1（即大于平均特征值）的主成分；几种准则应结合专业解释使用。例3只保留$c_1$，相对重构误差为$\lambda_2/2=3.4\%$。
\section{资料是否适合作主成分分析}
\subsection{Bartlett球形检验}
零假设为总体相关阵是单位阵（球形假设）。若$P<\alpha$，则认为指标间有一定相关性。
\begin{equation}
\chi^2=-\left(n-1-\frac{2p+5}{6}\right)\log|R|\ \dot\sim\ \chi^2\!\left(\frac{p(p-1)}{2}\right).
\end{equation}
\Note{$n$是总例数，$p$是变量个数。检验要求各观测独立抽样、变量服从多元正态分布，$\chi^2$分布是大样本近似。}
拒绝球形假设只说明变量间存在相关，并不说明相关强到足以有效降维：样本量大时，很弱的相关也会显著。降维效果仍要看特征值和累计贡献率。例2（表\ref{tab:hwv}）的$\chi^2=41.60$，$df=3$，$P<0.001$。
\begin{figure}[!htbp]\centering
\begin{tikzpicture}\begin{axis}[width=9cm,height=5.7cm,xlabel={$|R|$},ylabel={$\chi^2$},xmin=0,xmax=1,ymin=0,ymax=235,xtick={0,.2,.4,.6,.8,1},ytick={0,50,100,150,200}]
\addplot[Forest,line width=1pt] coordinates {(0.003003,144.715447) (0.012012,87.804878) (0.021021,76.016260) (0.030030,69.105691) (0.039039,63.414634) (0.048048,59.756098) (0.057057,56.504065) (0.066066,53.252033) (0.075075,50.813008) (0.084084,48.373984) (0.093093,46.747967) (0.102102,45.121951) (0.111111,43.089431) (0.120120,41.869919) (0.129129,40.243902) (0.138138,39.024390) (0.147147,37.804878) (0.156156,36.585366) (0.165165,35.772358) (0.174174,34.552846) (0.183183,33.739837) (0.192192,32.520325) (0.201201,31.707317) (0.210210,30.894309) (0.219219,30.081301) (0.228228,29.268293) (0.237237,28.455285) (0.246246,28.048780) (0.255255,27.235772) (0.264264,26.829268) (0.273273,26.016260) (0.282282,25.203252) (0.291291,24.796748) (0.300300,24.390244) (0.309309,23.577236) (0.318318,22.764228) (0.327327,22.764228) (0.336336,21.951220) (0.345345,21.544715) (0.354354,21.138211) (0.363363,20.325203) (0.372372,20.325203) (0.381381,19.512195) (0.390390,19.105691) (0.399399,18.699187) (0.408408,18.292683) (0.417417,17.886179) (0.426426,17.479675) (0.435435,17.073171) (0.444444,16.666667) (0.453453,16.260163) (0.462462,15.853659) (0.471471,15.447154) (0.480480,15.447154) (0.489489,14.634146) (0.498498,14.634146) (0.507508,13.821138) (0.516517,13.821138) (0.525526,13.414634) (0.534535,13.008130) (0.543544,13.008130) (0.552553,12.195122) (0.561562,12.195122) (0.570571,12.195122) (0.579580,11.382114) (0.588589,11.382114) (0.597598,10.975610) (0.606607,10.569106) (0.615616,10.569106) (0.624625,10.162602) (0.633634,9.756098) (0.642643,9.756098) (0.651652,9.349593) (0.660661,8.943089) (0.669670,8.943089) (0.678679,8.536585) (0.687688,8.130081) (0.696697,8.130081) (0.705706,8.130081) (0.714715,7.317073) (0.723724,7.317073) (0.732733,7.317073) (0.741742,6.504065) (0.750751,6.504065) (0.759760,6.504065) (0.768769,6.504065) (0.777778,5.691057) (0.786787,5.691057) (0.795796,5.691057) (0.804805,5.284553) (0.813814,4.878049) (0.822823,4.878049) (0.831832,4.878049) (0.840841,4.471545) (0.849850,4.065041) (0.858859,4.065041) (0.867868,4.065041) (0.876877,3.658537) (0.885886,3.252033) (0.894895,3.252033) (0.903904,3.252033) (0.912913,3.252033) (0.921922,2.439024) (0.930931,2.439024) (0.939940,2.439024) (0.948949,2.439024) (0.957958,1.626016) (0.966967,1.626016) (0.975976,1.626016) (0.984985,1.626016) (0.993994,1.626016)};
\end{axis}\end{tikzpicture}
\caption{Bartlett球形检验：卡方值随相关阵行列式的变化（示意图）}
\end{figure}\FloatBarrier

\subsection{KMO统计量}
\begin{equation}
\mathrm{KMO}=\frac{\sum_{i\ne j}r_{ij}^{2}}{\sum_{i\ne j}r_{ij}^{2}+\sum_{i\ne j}p_{ij}^{2}}.
\end{equation}
其中$p_{ij}$为偏相关系数。令$C=R^{-1}$，$C_{ij}$为逆矩阵$C$的元素，则
\[
p_{ij}=-\frac{C_{ij}}{\sqrt{C_{ii}C_{jj}}}.
\]
例2按表\ref{tab:hwv}复算，KMO$=0.695$（原例记为0.69）。
\begin{center}
\begin{tabular}{ll}\toprule
{\bfseries KMO范围}&{\bfseries 适合性评价}\\\midrule
$\mathrm{KMO}\geqslant0.9$&极好（Excellent）\\
$0.8\leqslant\mathrm{KMO}<0.9$&好（Good）\\
$0.7\leqslant\mathrm{KMO}<0.8$&中等（Fair）\\
$0.5\leqslant\mathrm{KMO}<0.7$&勉强接受（Poor）\\
$\mathrm{KMO}<0.5$&不可接受（Unacceptable）\\\bottomrule
\end{tabular}
\end{center}
例2的两个统计量可直接按公式计算，无需调用其他程序包：
\par\noindent\begin{rcode}
# h、w、v：例2表中29名男童的身高、体重、肺活量，按例1的方式录入
X <- cbind(h,w,v); R <- cor(X); n <- nrow(X); p <- ncol(X)
chi <- -(n-1-(2*p+5)/6)*log(det(R))
c(chi=chi, df=p*(p-1)/2, P=pchisq(chi,p*(p-1)/2,lower.tail=FALSE))
C <- solve(R); P <- -C/sqrt(outer(diag(C),diag(C)))   # 偏相关阵
diag(P) <- 0; R0 <- R; diag(R0) <- 0
sum(R0^2)/(sum(R0^2)+sum(P^2))                        # KMO
\end{rcode}
\begin{routput}
         chi           df            P 
4.159713e+01 3.000000e+00 4.884930e-09 
[1] 0.6945357
\end{routput}
KMO需与Bartlett球形检验结合使用。$\mathrm{KMO}\geqslant0.7$且Bartlett检验$P<0.05$时，数据适合因子分析。
\Note{KMO比较简单相关与偏相关：若变量间的相关主要来自共同因子，控制其他变量后偏相关应很小，KMO接近1。因此上表的经验阈值主要用于评价共同因子结构，是因子分析的参考指标。PCA只是对协方差结构作正交变换，不以存在公因子为前提，KMO不宜作为实施PCA的必要条件；PCA是否有用，应直接看前几个主成分的累计贡献率及其可解释性。}
\section{主成分应用}
\subsection{综合评价}
例3中，基于相关阵提取的第一主成分结构符合专业意义。假设成绩服从正态分布，可将主成分得分转化为百分位数（正态分布累积概率），获得较直观的解释。参照人群中$c_1\sim N(0,\lambda_1)$，故个体的百分位为
\[
P=100\,\Phi\!\left(\frac{c_1}{\sqrt{\lambda_1}}\right),
\]
即参照人群中得分低于该个体的百分比。这里的参照人群是估计均数、标准差和特征向量所用的人群，本例即这18名学生；换一个参照人群，百分位就要重新计算。特征向量的符号是任意的，应先把方向取成“得分越高越好”（若软件给出的$e_1$各分量为负，就用$-c_1$）。例3中$\lambda_1=1.932$：第1名学生（45，43）的$c_1=-2.85$，百分位为2.0；第9名（78，83）为0.33，百分位59.3；第17名（99，92）为1.59，百分位87.3。
\Note{方差最大准则只保证$c_1$最能拉开个体间的差距，并不保证它就是专业上合理的综合评价。还要检查各变量的系数是否同号、大小是否合理；若有变量方向相反（如“越小越好”的指标），应先统一方向再分析。}
\Example{例4\quad 空气污染加权指数}
广州市某两个月空气污染数据包括二氧化硫（SO$_2$）、氮氧化物（NO$_x$）和总悬浮颗粒物（TSP）三个原始浓度指标。各污染物分指数（iSO$_2$、iNO$_x$和iTSP）由浓度经非线性变换获得。api（Air Pollution Index）是当日三个分指数的最大值。

同样的分指数具有相同的医学意义和尺度，因此可以从协方差阵出发计算主成分，并构造空气污染加权指数\texttt{api.w}：$\operatorname{cor}(z_1,\mathrm{api.w})=1$。

只看相关系数等于1并不能确定这个指数：$z_1$的任何正线性变换$az_1+b$（$a>0$）与$z_1$的相关系数都等于1。课程程序（数据：\texttt{广州某2个月空气污染指数.sav}，$n=32$天）的具体做法是：对三个分指数的协方差阵求第一特征向量$e_1$（第一主成分贡献率76.4\%），取各分量为正的方向，再除以分量之和，使权重之和为1：
\[
a=\frac{e_1}{\sum_je_{1j}}=(0.322,\ 0.139,\ 0.540)\quad\text{对应（iTSP，iSO$_2$，iNO$_x$）},
\]
\[
\mathrm{api.w}=0.322\,\mathrm{iTSP}+0.139\,\mathrm{iSO_2}+0.540\,\mathrm{iNO_x}.
\]
也就是不中心化的加权平均，与分指数处于同一尺度。它与第一主成分得分$z_1=e_1\trans(x-\bar x)$的关系为$\mathrm{api.w}=z_1/\sum_je_{1j}+\bar x\trans a=0.643\,z_1+138.2$，即先缩放、再平移。iNO$_x$的权重最大，是因为它的方差最大；这是协方差阵分析按方差赋权的结果，并不表示它的健康效应最大。这个加权指数有无专业上的合理性，仍需结合污染物的健康效应来解释。
\begin{figure}[!htbp]\centering
\begin{tikzpicture}\begin{axis}[width=.95\textwidth,height=8.2cm,xlabel={day},ylabel={AQI},xmin=300,xmax=331,ymin=20,ymax=310,xtick={300,305,310,315,320,325,330},legend columns=3,legend style={at={(.5,1.1)},anchor=south,column sep=10pt}]\addplot[Forest,line width=1.25pt] coordinates {(300,107.26127) (301,145.27864) (302,141.22945) (303,157.21937) (304,81.621569) (305,90.085873) (306,86.178508) (307,98.646724) (308,78.291902) (309,72.049023) (310,104.56766) (311,84.434957) (312,86.18886) (313,213.97022) (314,152.9089) (315,147.39466) (316,157.99634) (317,154.27618) (318,74.077504) (319,128.10968) (320,169.44181) (321,200.48644) (322,174.34601) (323,148.40829) (324,84.311319) (325,85.25623) (326,110.46081) (327,88.88253) (328,117.71726) (329,171.52345) (330,197.32004) (331,179.32267)};\addlegendentry{api.w}
\addplot[Brick,densely dashed,line width=1pt] coordinates {(300,148) (301,217) (302,210) (303,212) (304,99) (305,107) (306,98) (307,120) (308,89) (309,84) (310,128) (311,101) (312,97) (313,259) (314,215) (315,207) (316,209) (317,194) (318,91) (319,181) (320,227) (321,261) (322,252) (323,211) (324,97) (325,102) (326,143) (327,133) (328,155) (329,253) (330,299) (331,270)};\addlegendentry{api}
\addplot[only marks,mark=*,mark size=1.7pt,Plum] coordinates {(300,148) (301,217) (302,210) (303,212) (304,99) (305,98) (306,98) (307,120) (308,89) (309,84) (310,128) (311,101) (312,97) (313,230) (314,215) (315,207) (316,209) (317,194) (318,91) (319,181) (320,227) (321,261) (322,252) (323,211) (324,97) (325,102) (326,143) (327,80) (328,155) (329,253) (330,299) (331,270)};\addlegendentry{iNO$_x$}
\addplot[only marks,mark=triangle,mark size=2.4pt,Ochre,line width=.65pt] coordinates {(300,80) (301,76) (302,73) (303,81) (304,73) (305,107) (306,94) (307,95) (308,82) (309,74) (310,103) (311,87) (312,97) (313,259) (314,91) (315,89) (316,116) (317,127) (318,69) (319,83) (320,130) (321,163) (322,93) (323,82) (324,86) (325,84) (326,91) (327,133) (328,93) (329,86) (330,80) (331,77)};\addlegendentry{iTSP}
\addplot[only marks,mark=+,mark size=2.4pt,Muted,line width=.65pt] coordinates {(300,12) (301,27) (302,32) (303,121) (304,34) (305,20) (306,22) (307,24) (308,28) (309,21) (310,17) (311,14) (312,19) (313,47) (314,55) (315,51) (316,57) (317,63) (318,20) (319,27) (320,37) (321,52) (322,61) (323,59) (324,31) (325,23) (326,29) (327,21) (328,30) (329,53) (330,74) (331,64)};\addlegendentry{iSO$_2$}\end{axis}\end{tikzpicture}
\caption{空气污染指数、加权指数与三个污染物分指数}
\end{figure}\FloatBarrier

\subsection{主成分回归}
22例胎儿受精龄（$Y$，周）与胎儿外型测量指标如下。试求由$X_1$、$X_2$、$X_3$推算$Y$的回归方程。
\begin{center}\begin{tabular}{lrr}\toprule
{\bfseries 指标}&{\bfseries 均数}&{\bfseries 标准差}\\\midrule
身高（$X_1$，cm）&33.05&9.71\\
头围（$X_2$，cm）&23.26&6.86\\
体重（$X_3$，g）&936.9&690.3\\\bottomrule
\end{tabular}\end{center}
因自变量间高度共线性，导致$b_2$的符号异常，为负值：
\[
\widehat Y=11.01+1.69X_1-2.16X_2+0.007X_3.
\]
\Note{主成分回归舍弃了部分主成分方向，一般是有偏估计；岭回归也是有偏估计。偏倚大小见下文推导。}
主成分回归步骤：基于相关阵提取1至2个主成分$c_1$和$c_2$，与$y$作回归（\texttt{y\string~c1+c2}），再变形为\texttt{y\string~x1+x2+x3}。
\begin{align*}
\widehat Y&=23.73+3.88C_1+3.10C_2,\\
C_1&=0.58\frac{X_1-\bar X_1}{S_1}+0.58\frac{X_2-\bar X_2}{S_2}+0.57\frac{X_3-\bar X_3}{S_3},\\
C_2&=-0.42\frac{X_1-\bar X_1}{S_1}-0.39\frac{X_2-\bar X_2}{S_2}+0.82\frac{X_3-\bar X_3}{S_3}.
\end{align*}
\subsubsection{由主成分回归还原原始变量的回归式}
\begin{align*}
\widehat Y&=b_0^c+b_1^cC_1+b_2^cC_2,\\
C_i&=a_{i1}x_1+a_{i2}x_2+a_{i3}x_3\\
&=\frac{a_{i1}(X_1-\bar X_1)}{S_1}+\frac{a_{i2}(X_2-\bar X_2)}{S_2}+\frac{a_{i3}(X_3-\bar X_3)}{S_3},\quad i=1,2.
\end{align*}
$b_i^c$为主成分回归系数，$a_{ij}$为组合系数（特征向量）。解得
\begin{align}
\widehat Y&=b_0+\frac{b_1^ca_{11}+b_2^ca_{21}}{S_1}X_1
+\frac{b_1^ca_{12}+b_2^ca_{22}}{S_2}X_2
+\frac{b_1^ca_{13}+b_2^ca_{23}}{S_3}X_3,\\
b_0&=b_0^c-\left[b_1^c\left(\frac{a_{11}\bar X_1}{S_1}+\frac{a_{12}\bar X_2}{S_2}+\frac{a_{13}\bar X_3}{S_3}\right)
+b_2^c\left(\frac{a_{21}\bar X_1}{S_1}+\frac{a_{22}\bar X_2}{S_2}+\frac{a_{23}\bar X_3}{S_3}\right)\right]\notag\\
&=b_0^c-(b_1\bar X_1+b_2\bar X_2+b_3\bar X_3).\notag
\end{align}
\subsubsection{分析结果}
\[
\widehat Y_{\mathrm{Prin}}(2)=10.44+0.10X_1+0.15X_2+0.007X_3.
\]
附：本例岭回归分析结果
\[
\widehat Y_{\mathrm{Ridge}}(0.1)=8.12+0.21X_1+0.19X_2+0.005X_3.
\]
\subsubsection{截断主成分回归的偏倚与方差}
设$X^*$为$n\times p$标准化自变量阵，$y=\beta_0\mathbf 1+X^*\beta+\varepsilon$，$\E\varepsilon=0$，$\cov(\varepsilon)=\sigma^2I$。由$X^{*\mathsf T}X^*=(n-1)E\Lambda E\trans$，主成分得分$Z=X^*E$各列两两正交，$Z\trans Z=(n-1)\Lambda$。记$\gamma=E\trans\beta$，即真回归系数在各主成分方向上的分量，则$X^*\beta=Z\gamma$。由于各列正交，$y$对$Z$回归时各系数互不影响：
\[
\widehat\gamma_j=\frac{z_j\trans y}{(n-1)\lambda_j},\qquad \E\widehat\gamma_j=\gamma_j,\qquad\var(\widehat\gamma_j)=\frac{\sigma^2}{(n-1)\lambda_j}.
\]
只保留前$k$个主成分时，$\widehat\beta_{(k)}=\sum_{j\leqslant k}\widehat\gamma_je_j$，于是
\[
\E\widehat\beta_{(k)}-\beta=-\sum_{j>k}\gamma_je_j,\qquad
\cov\bigl(\widehat\beta_{(k)}\bigr)=\frac{\sigma^2}{n-1}\sum_{j\leqslant k}\frac{e_je_j\trans}{\lambda_j}.
\]
偏倚只取决于真回归系数在被舍弃方向上的分量$\gamma_j$（$j>k$）：若$\beta$几乎与这些方向正交，偏倚很小。被舍弃的小$\lambda_j$方向本会贡献很大的方差$\sigma^2/[(n-1)\lambda_j]$，截断就是用偏倚换方差。保留全部非零主成分（$k=\operatorname{rank}(X^*)$）时，$Z_k$与$X^*$张成同一列空间，主成分回归恢复最小二乘拟合；本例$k=3$时还原得到的正是上面$b_2<0$的最小二乘方程。
\paragraph{原例续算}
本例（数据：\texttt{princeple reg exp.sav}）相关阵的特征值为2.926、0.0714、0.0025，前两个主成分的累计贡献率达99.9\%。逐个加入主成分时，$\widehat\gamma_1=3.88$、$\widehat\gamma_2=3.10$保持不变，与上面的推导一致。第三个主成分方差只有0.0025，但$\widehat\gamma_3=22.16$（标准误5.11，$t=4.34$，$P<0.001$），即低方差方向仍含有关于$Y$的预测信息。自变量的累计贡献率只说明$X$的结构，不能代替对$Y$的预测评价。用留一法比较预测误差，每一折内都重新标准化、重新作PCA：
\begin{code}
dc \ROperator{<-} foreign\ROperator{::}\RFunction{read.spss}(\RString{"princeple reg exp.sav"}, to.data.frame \ROperator{=} \RKeyword{TRUE})
X \ROperator{<-} \RFunction{as.matrix}(dc[, \RFunction{c}(\RString{"x1"}, \RString{"x2"}, \RString{"x3"})]); y \ROperator{<-} dc\ROperator{\$}y
rmse \ROperator{<-} \RFunction{sapply}(\RNumber{1}\ROperator{:}\RNumber{3}, \RKeyword{function}(k) \{
  e \ROperator{<-} \RFunction{numeric}(\RFunction{nrow}(X))
  \RKeyword{for} (i \RKeyword{in} \RFunction{seq_len}(\RFunction{nrow}(X))) \{
    pc \ROperator{<-} \RFunction{prcomp}(X[\ROperator{-}i, ], scale. \ROperator{=} \RKeyword{TRUE})      \RComment{\# 训练折内标准化与PCA}
    Z \ROperator{<-} pc\ROperator{\$}x[, \RNumber{1}\ROperator{:}k, drop \ROperator{=} \RKeyword{FALSE}]
    fit \ROperator{<-} \RFunction{lm}(y[\ROperator{-}i] \ROperator{~} Z)
    z0 \ROperator{<-} \RFunction{predict}(pc, X[i, , drop \ROperator{=} \RKeyword{FALSE}])[, \RNumber{1}\ROperator{:}k, drop \ROperator{=} \RKeyword{FALSE}]
    e[i] \ROperator{<-} y[i] \ROperator{-} \RFunction{cbind}(\RNumber{1}, z0) \ROperator{\%*\%} \RFunction{coef}(fit)
  \}
  \RFunction{sqrt}(\RFunction{mean}(e\ROperator{\textasciicircum{}}\RNumber{2}))
\})
rmse
\end{code}
\begin{center}\small\begin{tabular}{lccc}\toprule
{\bfseries 保留主成分数$k$}&1&2&3（即最小二乘）\\\midrule
留一法预测均方根误差（周）&1.867&1.732&1.483\\\bottomrule
\end{tabular}\end{center}
本例保留全部3个主成分时预测误差最小，PCR(2)虽然系数符号“合理”，预测反而较差。保留数应结合交叉验证的预测误差来选择；标准化和PCA都必须在训练折内估计，否则测试样品的信息会泄漏进模型。

最小二乘方程中$b_2=-2.16$（标准误0.54）是条件效应，表示身高、体重固定时头围每增加1 cm，$Y$的平均变化。$X_1$、$X_2$的方差膨胀因子约为200和209，条件数$\sqrt{2.926/0.0025}\approx34$，严重共线。这时“身高、体重不变而头围变化”的组合在数据中几乎不存在，负号的大小和方向都很不稳定，不宜直接作专业解释，应结合标准误和共线性诊断来看。

\begin{sikao}[主成分分析的几何概括]
主成分分析相当于建立一组新的坐标轴：第一轴沿数据分散程度最大的方向，第二轴取与第一轴垂直的方向中分散程度最大者，依此类推。仅保留前几个轴，即以尽量少的维度保留尽量多的方差。

本章的主要结论均可由此得到：方差最大对应特征向量；坐标轴相互垂直对应主成分互不相关；坐标轴旋转不改变总方差；分散程度按原始尺度度量，故结果依赖单位，单位不同时需先标准化；变量间无相关时数据在各方向上分散程度相同，无法降维。
\end{sikao}

\begin{chaptersummary}{本章小结：主成分分析}
研究目的&用少数互不相关的线性组合概括多个相关变量的方差，用于降维、综合评价、消除共线性\\
优化准则&在$a\trans a=1$且与前面各主成分不相关的约束下，使$\var(a\trans x)$最大；解为$S$或$R$的特征向量。等价于最小平方重构误差\\
主要条件&变量间有较强相关；单位不同时用相关阵；属于描述性方法，无分布假设（Bartlett检验另需多元正态）\\
结果解释&特征值即主成分方差，累计贡献率即重构出的方差比例；载荷$u_{i1}\sqrt{\lambda_1}$为变量与主成分的相关系数\\
常见误用&把“不相关”说成“独立”；单位不同仍用协方差阵；只凭累计贡献率决定主成分回归的保留数；把KMO当作PCA的必要条件；忽视特征向量符号的任意性\\
\end{chaptersummary}

\begin{fuxi}[复习题]
\begin{enumerate}
\item 对6个变量作主成分分析，分别写出基于相关阵与基于协方差阵时特征值之和。\\
\textit{答：相关阵为6；协方差阵为各变量方差之和$\tr(S)$。}
\item 解释“第一主成分的贡献率为70\%”的含义。\\
\textit{答：第一主成分的方差占总方差的70\%，即仅用第一主成分重构原数据（或标准化数据）时所保留的平方和比例。此处的“信息”专指方差。}
\item 基于相关阵的分析中，第一特征向量为$(0.58,0.58,0.57)$，$\lambda_1=2.926$，计算第一个变量与第一主成分的相关系数。\\
\textit{答：$0.58\times\sqrt{2.926}=0.99$。}
\item 说明变量间互不相关时不宜作主成分分析的原因。\\
\textit{答：此时相关阵接近单位阵，各特征值均接近1，每个主成分只代表一个原变量，无法降维。}
\item 说明主成分回归缓解共线性的原理及其代价。\\
\textit{答：舍弃特征值很小的方向，消除这些方向上被放大的方差；若真实系数在被舍弃方向上有分量，估计将产生偏倚，预测效果也可能下降。}
\end{enumerate}
\end{fuxi}
